% Auto-generated section: Why QA Reduces Training Compute
% This file is generated by qa_theory_export.rs. Do not edit by hand.

\section{Training Compute Reduction via Harmonic Optimization}

We compare classical stochastic gradient descent (SGD) against Harmonic Gradient Descent (HGD) on a family of quadratic harmonic surfaces.

For a representative configuration (dimension $d=\,16$, tolerance $\\tau=\,1.0e-10$, learning rate $\\eta=\,0.200$), we measured:

\begin{itemize}
  \item SGD steps to tolerance: $232.2$ (mean)
  \item HGD steps to tolerance: $111.2$ (mean)
  \item Step speedup: $2.09\times$
  \item Compute ratio (SGD/HGD): $1.25\times$ (proxy operations)
\end{itemize}

\paragraph{Harmonic Update Rule.}
Let $g$ denote the gradient and $\\Phi$ a harmonic gating functional combining mod-9 (digital-root) and mod-24 (phase) cycles.
HGD applies
\[
  w_{t+1} \,=\, w_t \, -\, \eta \, \big( \\Phi \\odot g \big),
\]
restricting updates to resonance-aligned directions. This suppresses high-frequency jitter and collapses plateau phases.

\paragraph{Complexity.}
HGD’s gating is symbolic and local; it avoids matrix inversions while attenuating destructive curvature. Integer-harmonic implementations admit further bit-width reductions.

\paragraph{Conclusion.}
Across synthetic tasks, HGD achieves fewer steps and lower proxy compute than SGD, consistent with the harmonic constraints in QA.

\paragraph{Sweep Robustness.}
On a grid of benchmark settings (varying dimension and learning rate), QA-HGD won in $1$ out of $1$ configurations, with mean step speedup $2.09\times$ and mean compute ratio $1.25\times$.
